\chapter{Getting started with Julia}
\label{ch:getting-started}

\begin{quote}
\emph{``We want a language that's open source, with a liberal license. We want the speed of C with the dynamism of Ruby. We want homoiconicity, true macros, and a language that is as usable for general programming as Python.''}
--- Jeff Bezanson, Stefan Karpinski, Viral B. Shah, Alan Edelman (2012)
\end{quote}

% ============================================================
\section{Why Julia?}

Julia was created to solve the \emph{two-language problem}: scientists prototype in a high-level language (Python, Matlab, R) and then rewrite performance-critical code in C or Fortran. Julia eliminates this step. Code that reads like Python runs at speeds comparable to C.

Key features that make Julia distinctive:

\begin{enumerate}
  \item \textbf{Just-in-time (JIT) compilation.} Julia compiles functions to native machine code using LLVM the first time they are called. Subsequent calls are fast.
  \item \textbf{Multiple dispatch.} Functions can have many methods, selected based on the types of \emph{all} arguments. This is Julia's central design principle.
  \item \textbf{Rich type system.} Abstract types, parametric types, and union types let you write generic, reusable code.
  \item \textbf{Metaprogramming.} Julia code is represented as Julia data structures (expressions), enabling powerful macros.
  \item \textbf{Package ecosystem.} Over 10,000 registered packages for scientific computing, data science, optimization, and machine learning.
\end{enumerate}

\begin{juliatip}
Julia uses 1-based indexing, like Matlab and R, not 0-based like Python and C. The first element of an array \texttt{v} is \texttt{v[1]}.
\end{juliatip}

% ============================================================
\section{Installing Julia}

\subsection{Using juliaup (recommended)}

\begin{minted}{bash}
# macOS / Linux
curl -fsSL https://install.julialang.org | sh

# Windows (PowerShell)
winget install julia -s msstore
\end{minted}

Then verify:
\begin{minted}{bash}
julia --version
# julia version 1.10.x
\end{minted}

\subsection{VS Code integration}

Install the Julia extension in VS Code. It provides:
\begin{itemize}
  \item Syntax highlighting and autocompletion
  \item Integrated REPL (\texttt{Alt+J, Alt+O} to open)
  \item Plot pane for inline visualization
  \item Debugger and profiler integration
\end{itemize}

% ============================================================
\section{The Julia REPL}

Launch Julia by typing \texttt{julia} in your terminal. The REPL (Read-Eval-Print Loop) has four modes:

\begin{minted}{julia}
# Julian mode (default) — execute Julia code
julia> 2 + 3
5

julia> sqrt(144)
12.0
\end{minted}

\begin{minted}{julia}
# Package mode — press ]
(@v1.10) pkg> add DataFrames

# Help mode — press ?
help?> sqrt
  sqrt(x)
  Return the square root of x.

# Shell mode — press ;
shell> ls
\end{minted}

\begin{juliatip}
Press \texttt{Backspace} on an empty line to return to Julian mode from any special mode. Press \texttt{Tab} for autocompletion — it works on function names, file paths, and even Unicode characters.
\end{juliatip}

% ============================================================
\section{Package management with Pkg}

Julia's built-in package manager is powerful and reproducible:

\begin{minted}{julia}
using Pkg

# Install a package
Pkg.add("DataFrames")

# Install multiple packages at once
Pkg.add(["CSV", "Plots", "BenchmarkTools"])

# Update all packages
Pkg.update()

# Check installed packages
Pkg.status()

# Create a project environment (recommended for reproducibility)
Pkg.activate("MyProject")
Pkg.add("DataFrames")
# This creates Project.toml and Manifest.toml
\end{minted}

\begin{warningbox}
Always use project environments for serious work. The default global environment (\texttt{@v1.10}) can develop dependency conflicts if you install too many packages. Run \texttt{Pkg.activate(".")} in your project directory to create a local environment.
\end{warningbox}

% ============================================================
\section{Your first Julia program}

\begin{minted}{julia}
# hello.jl — Your first Julia script

# Variables
name = "Julia"
year = 2012
age = 2026 - year

println("Hello from $name!")
println("Julia is $age years old and faster than ever.")

# A simple computation
function fibonacci(n)
    n <= 1 && return n
    a, b = 0, 1
    for _ in 2:n
        a, b = b, a + b
    end
    return b
end

# Print the first 15 Fibonacci numbers
for i in 0:14
    println("F($i) = $(fibonacci(i))")
end
\end{minted}

Run this script from the terminal:
\begin{minted}{bash}
julia hello.jl
\end{minted}

Or from the REPL:
\begin{minted}{julia}
julia> include("hello.jl")
\end{minted}

% ============================================================
\section{Pluto notebooks}

Pluto is Julia's reactive notebook environment. Unlike Jupyter, Pluto notebooks are:
\begin{itemize}
  \item \textbf{Reactive:} changing one cell automatically updates all cells that depend on it.
  \item \textbf{Reproducible:} cells execute in dependency order, not in the order you click.
  \item \textbf{Pure Julia files:} saved as \texttt{.jl} files, not JSON. Easy to version-control with Git.
\end{itemize}

\begin{minted}{julia}
# Install and launch Pluto
using Pkg
Pkg.add("Pluto")

using Pluto
Pluto.run()   # opens a browser tab at localhost:1234
\end{minted}

Inside a Pluto notebook, each cell contains one expression or \texttt{begin...end} block:

\begin{minted}{julia}
# Cell 1
x = 42

# Cell 2 — automatically re-runs when x changes
y = x^2 + 1

# Cell 3
md"The value of y is **$(y)**"
\end{minted}

\begin{juliatip}
Pluto enforces one assignment per variable per notebook. This prevents hidden state bugs. If you need multiple statements in one cell, wrap them in a \texttt{begin...end} block.
\end{juliatip}

% ============================================================
\section{IJulia and Jupyter}

If you prefer Jupyter notebooks (e.g., for sharing with Python colleagues):

\begin{minted}{julia}
using Pkg
Pkg.add("IJulia")

using IJulia
notebook()   # opens Jupyter in the browser with Julia kernel
\end{minted}

% ============================================================
\section{Unicode in Julia}

Julia supports Unicode variable names and operators. Type them using LaTeX-like tab completion:

\begin{minted}{julia}
# Type \alpha then press Tab to get α
α = 0.05
β = 1 - α

# Type \sqrt then Tab to get √
√2 = √(2)    # equivalent to sqrt(2)

# Greek letters in functions
function Δ(a, b, c)
    return b^2 - 4a*c
end

# Matrix operations with Unicode
A = [1 2; 3 4]
x⃗ = [5, 6]           # \vec then Tab after x
result = A * x⃗
\end{minted}

\begin{exercise}
\begin{enumerate}
  \item Install Julia 1.10 on your machine using \texttt{juliaup}. Verify by running \texttt{julia --version}.
  \item In the REPL, compute $\sqrt{2^{10} + 3^{10}}$ and $\sum_{k=1}^{100} \frac{1}{k^2}$.
  \item Install Pluto.jl and create a notebook with three cells: (a) define $n = 20$, (b) compute $n!$ using \texttt{factorial(n)}, (c) display the result in a Markdown cell. Change $n$ and observe reactivity.
  \item Create a project environment called \texttt{JuliaCourse}, add the packages \texttt{DataFrames}, \texttt{CSV}, and \texttt{Plots}. Check the \texttt{Project.toml} file.
  \item Write a function \texttt{collatz(n)} that returns the number of steps to reach 1 in the Collatz sequence. Test it on $n = 27$.
\end{enumerate}
\end{exercise}

% ============================================================
\section{Chapter summary}

\begin{itemize}
  \item Julia solves the two-language problem: high-level syntax with compiled speed.
  \item The REPL has four modes: Julian, Package (\texttt{]}), Help (\texttt{?}), and Shell (\texttt{;}).
  \item Use \texttt{Pkg.activate(".")} to create reproducible project environments.
  \item Pluto.jl provides reactive notebooks saved as plain Julia files.
  \item IJulia integrates Julia with the Jupyter ecosystem.
  \item Julia supports Unicode natively --- use Greek letters and math symbols freely.
\end{itemize}
